\section{Results and Discussion}
\subsection{Results and Analysis}

\subsubsection{Barren-Plateau Scaling}
\textbf{[PENDING re-measurement -- see note below.]} The gradient variance measured on the AI4I 2020 predictive-maintenance dataset (the project's primary dataset; see Section~4.2) was previously reported as dropping by a factor of 57.3$\times$ from 6 to 12 qubits, close to the theoretical $2^{-n}$ expectation. That specific comparison is now out of scope: the project's validated qubit range on available hardware is $n\le8$ (Section~1.6), so the scaling claim needs to be re-measured over 6 to 8 qubits, where the theoretical $2^{-n}$ expectation is a much smaller $4\times$ drop rather than $64\times$ -- a genuinely weaker signal that has not yet been confirmed empirically at this narrower range. \Cref{tab:ai4i} below is pending the same re-scoped re-run, together with the SMOTE-related data-pipeline correction noted after it.

\begin{table}[htbp]
\centering
\caption{Scheduler comparison on AI4I (n = 8, local cost, mean over 3 seeds) -- pending re-run.}
\label{tab:ai4i}
\begin{tabular}{|p{3cm}|p{2cm}|p{2.3cm}|p{2cm}|p{2.5cm}|}
\hline
\textbf{Scheduler} & \textbf{Type} & \textbf{Accuracy} & \textbf{F1} & \textbf{Depth} \\
\hline
Rule & static & -- & -- & -- \\
Greedy & static & -- & -- & -- \\
PPO & learned & -- & -- & -- \\
DQN & learned & -- & -- & -- \\
\hline
\end{tabular}
\end{table}

\textbf{Data-pipeline and scope note.} This table previously reported $n=12$ results generated before two subsequent changes: (i) the qubit range validated on available hardware was narrowed to $n\le8$ (Section~1.6, RAM/runtime cost of PennyLane statevector simulation), which alone requires this benchmark to move from $n=12$ to $n=8$; and (ii) the AI4I loader was changed from forcing a 339/339 undersampled split to preserving the natural $\approx$3.4\% class imbalance and correcting it with SMOTE on the training split only (Section~4.2, AI4I 2020 Predictive Maintenance). Both changes require the primary benchmark to be re-run at $n=8$ against the corrected loader before this table can be filled in; the old $n=12$ numbers have been removed rather than left in place mislabelled as $n=8$.

\subsubsection{Primary AI4I Benchmark}
\textbf{[PENDING -- to be written once \Cref{tab:ai4i} is re-run at $n=8$; see note above.]}

\subsubsection{Real Deployment versus a Synthetic Comparison}
Beyond the primary AI4I evaluation above, a second block of results combines one real deployment on an independent dataset with a synthetic comparison at the same qubit count. \Cref{tab:heart} is the real, physical anchor: the policy deployed on the actual PennyLane statevector circuit for Heart Disease at 8 qubits under a local cost. \Cref{tab:sim8g} varies only the cost type (global instead of local) at the same 8-qubit count on synthetic data, isolating the effect of cost type on trainability and scheduler behaviour while holding qubit count fixed; because the cost type and dataset differ between these two tables, they should be read as separate regimes rather than as a single matched comparison.

\begin{table}[htbp]
\centering
\caption{Scheduler comparison -- PennyLane real deployment (Heart Disease, n = 8 qubits, local cost, single seed).}
\label{tab:heart}
\begin{tabular}{|p{3cm}|p{2cm}|p{2.3cm}|p{2cm}|p{2.5cm}|}
\hline
\textbf{Scheduler} & \textbf{Type} & \textbf{Accuracy} & \textbf{F1} & \textbf{Depth} \\
\hline
Rule & static & 0.518 & 0.608 & 4 \\
Greedy & static & 0.807 & 0.795 & 16 \\
PPO & learned & 0.735 & 0.703 & 2 \\
DQN & learned & 0.554 & 0.667 & 2 \\
\hline
\end{tabular}
\end{table}

\subsubsection{PennyLane Real Deployment (Heart Disease)}
On the real PennyLane deployment on Heart Disease at 8 qubits with a local cost, the greedy heuristic achieves the best accuracy (0.807) and F1 (0.795), growing the circuit to depth 16; the fixed-depth baselines at k = 8 and k = 16 perform almost identically (accuracy 0.795), showing that once the circuit is deep enough under a local cost the gradient remains trainable and extra depth adds little further benefit. PPO matches a shallow fixed-depth baseline (accuracy 0.735 at depth 2) without growing the circuit, while the hand-tuned rule and DQN both underperform, with DQN in particular collapsing to a near-chance ROC AUC (0.492). The ranking Greedy > PPO > Rule > DQN on this real, trainable deployment is a useful counterexample to any claim that one scheduler dominates universally.

\begin{table}[htbp]
\centering
\caption{Scheduler comparison -- simulated deployment (synthetic data, n = 8 qubits, global cost, single seed).}
\label{tab:sim8g}
\begin{tabular}{|p{3cm}|p{2cm}|p{2.3cm}|p{2cm}|p{2.5cm}|}
\hline
\textbf{Scheduler} & \textbf{Type} & \textbf{Accuracy} & \textbf{F1} & \textbf{Depth} \\
\hline
Rule & static & 0.680 & 0.666 & 6 \\
Greedy & static & 0.908 & 0.889 & 16 \\
PPO & learned & 0.735 & 0.720 & 7 \\
DQN & learned & 0.680 & 0.666 & 6 \\
\hline
\end{tabular}
\end{table}

\subsubsection{Same-Qubit-Count Synthetic Comparison (Global Cost)}
At the same qubit count as the real deployment (8) but with a global cost, the synthetic run shows the same qualitative ranking as \Cref{tab:heart} -- greedy on top (0.908), PPO next (0.735) -- but with higher absolute accuracy across the board. DQN and the hand-tuned rule land on an identical accuracy (0.680) at the same depth (6), an exact tie that is a coincidence of this single seed rather than a general result. The qualitative agreement with the real \Cref{tab:heart} ranking is reassuring for the synthetic pipeline, but the different cost type (global vs local) means the gap in absolute accuracy should not be over-interpreted.

\subsubsection{Additional Single-Dataset Deployments (German Credit, Parkinson's, Pima)}
Beyond the core evaluation design above, the project also logged one real PennyLane deployment run (8 qubits, single seed) on each of three secondary datasets from \Cref{tab:datasets} -- German Credit, Parkinson's, and Pima -- using the same rule/greedy/PPO/DQN comparison procedure as \Cref{tab:heart}. These were exploratory checks rather than part of the core evaluation design (Section~1.4), and a separate run-configuration file was not retained alongside the logs, so the cost type and seed are inferred from the comparison pipeline's default configuration (local cost, seed 123) rather than confirmed from a recorded config; they are reported here, with that caveat, because the logs already exist and were otherwise undocumented.

\begin{table}[htbp]
\centering
\caption{Scheduler comparison -- PennyLane real deployment (German Credit, n = 8 qubits, single seed; cost type inferred as local from the pipeline default).}
\label{tab:german}
\begin{tabular}{|p{3cm}|p{2cm}|p{2.3cm}|p{2cm}|p{2.5cm}|}
\hline
\textbf{Scheduler} & \textbf{Type} & \textbf{Accuracy} & \textbf{F1} & \textbf{Depth} \\
\hline
Rule & static & 0.706 & 0.640 & 6 \\
Greedy & static & 0.667 & 0.706 & 3 \\
PPO & learned & 0.722 & 0.734 & 6 \\
DQN & learned & 0.628 & 0.511 & 1 \\
\hline
\end{tabular}
\end{table}

On German Credit, PPO attains the best accuracy (0.722) and F1 (0.734) of the four schedulers, narrowly ahead of the hand-tuned rule (0.706 accuracy) at the same committed depth (6); greedy stops early at depth 3, and DQN collapses to depth 1 with the weakest F1 (0.511). This is the one regime in this report where PPO outright wins -- a useful reminder that the ranking between the rule and PPO is not fixed across datasets, to be checked against the re-run AI4I benchmark (\Cref{tab:ai4i}) once available.

\begin{table}[htbp]
\centering
\caption{Scheduler comparison -- PennyLane real deployment (Parkinson's, n = 8 qubits, single seed; cost type inferred as local from the pipeline default).}
\label{tab:parkinsons}
\begin{tabular}{|p{3cm}|p{2cm}|p{2.3cm}|p{2cm}|p{2.5cm}|}
\hline
\textbf{Scheduler} & \textbf{Type} & \textbf{Accuracy} & \textbf{F1} & \textbf{Depth} \\
\hline
Rule & static & 0.899 & 0.892 & 6 \\
Greedy & static & 0.843 & 0.851 & 8 \\
PPO & learned & 0.730 & 0.745 & 7 \\
DQN & learned & 0.640 & 0.515 & 2 \\
\hline
\end{tabular}
\end{table}

On Parkinson's, the hand-tuned rule is clearly the strongest scheduler (0.899 accuracy, 0.892 F1), with greedy second; both static schedulers beat PPO and DQN by a wide margin. DQN is notably worse than its own 0.719 starting accuracy (\texttt{acc\_start} in the underlying log), meaning the agent's depth decisions actively hurt this small, 96-sample dataset rather than merely failing to help -- another instance of the DQN instability already seen on \Cref{tab:heart}.

\begin{table}[htbp]
\centering
\caption{Scheduler comparison -- PennyLane real deployment (Pima, n = 8 qubits, single seed; cost type inferred as local from the pipeline default).}
\label{tab:pima}
\begin{tabular}{|p{3cm}|p{2cm}|p{2.3cm}|p{2cm}|p{2.5cm}|}
\hline
\textbf{Scheduler} & \textbf{Type} & \textbf{Accuracy} & \textbf{F1} & \textbf{Depth} \\
\hline
Rule & static & 0.689 & 0.680 & 6 \\
Greedy & static & 0.714 & 0.742 & 8 \\
PPO & learned & 0.578 & 0.558 & 2 \\
DQN & learned & 0.708 & 0.704 & 2 \\
\hline
\end{tabular}
\end{table}

On Pima, greedy is best (0.714 accuracy), with DQN a close second (0.708) despite committing to a much shallower depth (2 versus greedy's 8); PPO is the weakest scheduler here and, as on Parkinson's, ends below its own starting accuracy (0.590), indicating an unproductive depth change rather than a neutral one. Across all three supplementary datasets, no single scheduler wins twice: PPO wins on German Credit, the rule wins on Parkinson's, and greedy wins on Pima. This adds three more data points to the same conclusion the core evaluation reaches -- that scheduler ranking is regime- and dataset-dependent -- while underscoring, precisely because these three runs are single-seed and only single-dataset checks, that they should be read as corroborating evidence rather than as an independent, equally-weighted evaluation alongside \Cref{tab:ai4i,tab:heart,tab:sim8g}.

\subsection{Extended Fixed-Depth Sweep}

\subsubsection{Sweep Methodology}
\Cref{tab:ai4i,tab:heart,tab:sim8g} summarize each regime with only four schedulers (Rule, Greedy, PPO, DQN) for readability. The deployment pipeline, however, also runs a dense fixed-depth sweep (k = 2, 4, 6, 8, 12, 16) as a reference curve against which the adaptive schedulers can be judged: if a learned or hand-tuned scheduler settles on a depth d, this sweep tells us whether d was actually close to the accuracy-maximizing depth for that regime, or whether the scheduler stopped short (or overshot). This subsection reports the full sweep for the two single-seed experimental configurations behind \Cref{tab:heart,tab:sim8g}, so that those summary rankings can be checked against the underlying accuracy-vs-depth curve rather than taken on faith; the three-seed primary AI4I benchmark (\Cref{tab:ai4i}) is a separate, higher-cost run and was not additionally swept across depths.


\begin{table}[htbp]
\centering
\caption{Fixed-depth sweep -- PennyLane real deployment (Heart Disease, n = 8 qubits, local cost).}
\label{tab:sweep_heart}
\begin{tabular}{|p{2.4cm}|p{2.3cm}|p{2cm}|p{2.4cm}|}
\hline
\textbf{Depth k} & \textbf{Accuracy} & \textbf{F1} & \textbf{ROC AUC} \\
\hline
2 & 0.735 & 0.703 & 0.835 \\
4 & 0.759 & 0.756 & 0.845 \\
8 & 0.795 & 0.790 & 0.866 \\
16 & 0.795 & 0.790 & 0.876 \\
\hline
\end{tabular}
\end{table}

\subsubsection{Heart Disease Sweep (Real Deployment)}
On the real Heart Disease deployment, accuracy rises steeply from depth 2 to depth 4 (+0.024) and then flattens almost completely between depth 8 and depth 16 (+0.000 accuracy, +0.010 ROC AUC), consistent with the local-cost gradient-variance measurement in Section~6.1: once the circuit is deep enough for the local cost to keep gradients alive, additional depth buys almost nothing. This is exactly why Greedy (which grows to depth 16 by construction) and the k = 8 fixed baseline land within 0.4 percentage points of each other -- the accuracy-maximizing depth on this real circuit is reached well before the depth cap, and the fixed-depth curve is the evidence, not just the scheduler comparison in \Cref{tab:heart}.

\begin{table}[htbp]
\centering
\caption{Fixed-depth sweep -- simulated deployment (synthetic, n = 8 qubits, global cost).}
\label{tab:sweep_sim8g}
\begin{tabular}{|p{2.4cm}|p{2.3cm}|p{2cm}|p{2.4cm}|}
\hline
\textbf{Depth k} & \textbf{Accuracy} & \textbf{F1} & \textbf{ROC AUC} \\
\hline
2 & 0.254 & 0.249 & 0.627 \\
4 & 0.519 & 0.509 & 0.760 \\
6 & 0.680 & 0.666 & 0.840 \\
8 & 0.777 & 0.762 & 0.889 \\
12 & 0.873 & 0.855 & 0.936 \\
16 & 0.453 & 0.444 & 0.727 \\
\hline
\end{tabular}
\end{table}

\subsubsection{Synthetic 8-Qubit Global-Cost Sweep}
The synthetic q = 8, global-cost sweep is monotonically increasing from depth 2 through depth 12, gaining roughly 0.10--0.15 accuracy per doubling of depth, before dropping sharply at depth 16 (accuracy falls from 0.873 at k = 12 to 0.453 at k = 16). This non-monotonic tail is a single-seed artefact of the calibrated surrogate rather than a physical prediction -- it is precisely the kind of extrapolation risk flagged in Section~4.7.1, where the surrogate is described as reliable only within its calibrated regime. It also explains why Greedy, which happens to land on depth 16 in \Cref{tab:sim8g}, still reports a high accuracy (0.908): Greedy's own trajectory through the depth-vs-accuracy landscape does not pass through the same point as this independently swept curve, since the two runs commit layers under different stopping criteria. The practical lesson is that a single fixed-depth sweep on a single seed should not be read as a precise accuracy-vs-depth function, only as a rough guide to where the trainable region lies.

\subsection{Discussion}

\subsubsection{Answering the Central Research Question}
\textbf{[PENDING -- to be finalized once the AI4I n=8 benchmark is available; see note below.]} The central research question -- whether an RL agent can schedule circuit depth to keep a variational classifier trainable -- is answered with a qualified yes, though the strength of that answer depends heavily on regime rather than on the agent alone. On the real Heart Disease deployment (\Cref{tab:heart}) and its synthetic same-qubit counterpart (\Cref{tab:sim8g}), both under cost regimes where gradients stay alive, static heuristics -- greedy in particular -- match or beat the learned agents; PPO tracks a shallow fixed-depth baseline instead of demonstrating a clear advantage from adaptive growth. The clearest case for adaptive scheduling in this report is instead the supplementary German Credit deployment (\Cref{tab:german}, Section~6.1), where PPO outright wins on both accuracy and F1 -- though that result is single-seed and exploratory rather than part of the core design. This is broadly consistent with the reward analysis: because the loss-improvement weight dominates the gradient-variance weight (\Cref{tab:reward}), the agent primarily learns to optimize loss under a resource constraint rather than to actively avoid barren plateaus, so a genuine plateau-avoidance advantage is the exception here rather than the rule. \textbf{Note:} an earlier version of this section additionally reported a 20-qubit, local-cost synthetic regime in which PPO was the clear winner; that regime has been removed because it falls outside the $\le$8-qubit scope validated on available hardware (Section~1.6, RAM/runtime cost of PennyLane statevector simulation), and this paragraph should be revisited once the new AI4I n=8 benchmark (\Cref{tab:ai4i}) is available, since it is the most likely remaining source of a clear scheduler-value result.

\subsubsection{Alignment with Theory}
These findings align with the theory: local cost functions are known to weaken barren plateaus while introducing depth dependence, which is the condition under which the scheduler's depth control has the potential to matter. The honest conclusion is not that reinforcement learning is a universal remedy, but that any advantage it offers is regime-specific and modest at the qubit counts evaluated here -- and that a strong static baseline must always be reported alongside it.

\subsubsection{Limitations Encountered}
Limitations encountered during the project include the instability of single-seed deployment -- \Cref{tab:heart,tab:sim8g} above each report a single seed, so headline numbers such as greedy's 0.807 accuracy on Heart Disease (\Cref{tab:heart}) should be read as illustrative rather than statistically robust; \Cref{tab:ai4i} (AI4I) is the only three-seed result in this section, and even that is limited to a single dataset and qubit count -- together with the cost of statevector evaluation and the sensitivity of results to preprocessing choices, which motivated the standardized fit-on-train pipeline.

\subsection{Recommendations}

\subsubsection{Scheduling and Evaluation Practice}
\begin{itemize}
\item Always report a hand-tuned rule and a greedy heuristic alongside the learned scheduler; adaptive scheduling should be claimed only where it demonstrably separates from them.
\item Evaluate every headline number with at least three seeds and report mean $\pm$ standard deviation; treat single-seed runs as exploratory.
\end{itemize}

\subsubsection{Regime Selection and Simulator Fidelity}
\begin{itemize}
\item Target the local-cost regime when depth scheduling is the goal, since that is where the plateau -- and therefore the potential benefit -- is depth-dependent rather than fixed.
\item Re-calibrate the surrogate whenever the ansatz or preprocessing changes, to keep the training environment faithful to deployment.
\end{itemize}

